% Job posting: https://ca.linkedin.com/jobs/view/modiface-sr-ai-researcher-at-l-or%C3%A9al-4463883471
\documentclass[11pt]{article}
\usepackage[left=1in,top=1in,right=1in,bottom=1in]{geometry}
\usepackage[hidelinks]{hyperref}
\usepackage{setspace}
\usepackage[T1]{fontenc}
\usepackage{newtxtext,newtxmath}
\ifdefined\pdfgentounicode
  \input{glyphtounicode}
  \pdfgentounicode=1
\fi
\setstretch{1.1}
\pagenumbering{gobble}
\begin{document}
\pagestyle{empty}

\noindent
Nishal Stanislaus Silva, PhD \\
Toronto, ON \\
nishal.silva@hotmail.com \,|\, +1 613 200 2030 \\
nishal.xyz \,|\, linkedin.com/in/nishal-silva \,|\, github.com/ns2max

\vspace{1.5em}
\noindent
Dear Hiring Manager,

\vspace{1em}
\noindent
I am writing to apply for the Sr. AI Researcher role at ModiFace, L'Or\'eal's AR/AI beauty-tech division. During my postdoctoral research at the University of Trento, I published a real-time pattern-detection system running at 14 ms inference on a Raspberry Pi 4 -- F1 0.76, 31.4\% CPU, 74 times faster than the 1038 ms baseline it replaced (JAES 2026). Shipping ML under a hard, imperceptible-latency budget on constrained edge hardware is exactly the discipline that AR try-on and face tracking demand to feel natural and instantaneous to a consumer holding a phone: any perceptible lag breaks the illusion the product depends on.

\vspace{1em}
\noindent
My background bridges production computer vision and rigorous evaluation methodology. At MAS Holdings, I designed and deployed OpenCV-based computer-vision systems at industrial scale -- including an automated care-label quality-control system using template and feature matching with an explainable defect overlay, and a loom-side fabric-inspection system using a microscopic camera array -- work that required the same attention to real-world robustness and deployment constraints that consumer AR must meet outside the lab. My research since has centred on frozen-protocol benchmark and ablation design, quantifying accuracy, latency, and compute trade-offs with a written record of negative results -- the kind of rigor needed to evaluate perceptual quality trade-offs in a virtual try-on pipeline. I also bring direct real-time mixed-reality systems experience: at Trento, I built a pipeline synchronizing audio, BCI/EEG, and head/hand tracking across four performers simultaneously using Unity and Meta Quest 3 headsets, run live in two public performances.

\vspace{1em}
\noindent
I have not worked in beauty-tech or consumer AR specifically, and I want to be upfront about that. What I bring instead is the transferable core the role requires: real-time, on-device computer vision and inference engineered against hard latency budgets, and benchmark design rigorous enough to make quality trade-offs measurable rather than assumed. I am confident that core translates directly to ModiFace's product surface.

\vspace{1em}
\noindent
I am a Canadian Permanent Resident, based in Toronto, eligible to work without sponsorship, and available immediately. I would welcome the chance to discuss how my real-time computer vision and benchmarking background could contribute to ModiFace's work.

\vspace{1.5em}
\noindent
Sincerely, \\
Nishal Stanislaus Silva

\end{document}
